\documentclass[10pt,a4paper]{amsart}
\usepackage[margin=19mm]{geometry}
\usepackage{amsmath,amssymb,mathtools}
\usepackage[scr=rsfs,cal=euler]{mathalfa}
\usepackage{xfrac}
\usepackage[unicode,hidelinks,bookmarks=false]{hyperref}
\newcommand{\ie}{i.e.\ }
\newcommand{\cf}{cf.\ }
\newcommand{\resp}{respectively\ }
\newcommand{\Subsection}{\subsection}
\providecommand{\nobreakdash}{\nobreak\hbox{-}}
\setlength{\emergencystretch}{2em}
\multlinegap0pt
\usepackage{bm}
\usepackage{bbm}
\usepackage{dsfont}
\usepackage{xspace}
\usepackage{cancel}
\usepackage{mathtools}
\usepackage{amsthm}
\makeatletter
\@ifundefined{theorem}{\newtheorem{theorem}{Theorem}}{}
\@ifundefined{lemma}{\newtheorem{lemma}[theorem]{Lemma}}{}
\makeatother
\newcommand\bx{\boldsymbol{x}}
\newcommand\calP{{\mathcal{P}}}
\newcommand\calT{{\mathcal{T}}}
\newcommand{\bnabla}{\nabla_{\kern-.05cm\mesh}^{}}
\newcommand{\jump}[1]{[\![#1]\!]}                        % jump
\newcommand{\mean}[1]{\{\kern-1.1mm\{#1\}\kern-1.1mm\}}          
\newcommand{\mesh}{\calP}
\newcommand{\ndgp}[1]{| \kern -.25mm \|{#1}| \kern -.25mm \|_{\rm \mesh}}
\newcommand{\norm}[2]{\|{#1}\|_{{#2}}}
\newcommand{\ud}{\,\mathrm{d}}
\renewcommand{\k}{K}
\title[A nodally bound-preserving composite discontinuous Galerkin ]{A nodally bound-preserving composite discontinuous Galerkin method on polytopic meshes}
\author{Abdolreza Amiri, Gabriel R. Barrenechea, Emmanuil H. Georgoulis, Tristan Pryer}
\date{Source: arXiv:2510.02094v1 (CC BY 4.0)}
\begin{document}
\maketitle

\noindent\emph{Source and licence.} A nodally bound-preserving composite discontinuous Galerkin method on polytopic meshes. Version arXiv:2510.02094v1 (CC BY 4.0), distributed under the Creative Commons Attribution 4.0 International licence, as recorded in the arXiv licence metadata for this exact version and shown on the abstract page https://arxiv.org/abs/2510.02094v1. Full rights notice: licenses/arxiv-2510.02094-CC-BY-4.0.txt.\par\smallskip

\noindent\textit{Editorial scope.} Counted unit: the Lemma whose source label is \texttt{Lem:cont\_twoscale} in the article (source0.tex lines 902--921, source label \texttt{Lem:cont\_twoscale}), delivered together with the article's own complete original proof of it (source0.tex lines 923--1026, one \texttt{proof} environment of 104 lines). No mathematics is written, repaired or re-derived: statement and proof are the article's own text line for line. Required context retained uncounted: indef\_est\_important (lines 653--657); eq20 (lines 727--731); Lem:s (lines 780--789). Every definition, display and label of those lines is retained unchanged. Omitted from this package and not redistributed: 1--652, 658--726, 732--779, 790--901, 922--922, 1027--2052. Nothing inside a retained excerpt is omitted or altered: every retained body is delivered exactly as the article's lines are written, so no omission receipt of any kind accompanies this package. Citations: the retained excerpts carry no citation command at all, so no bibliography entry of the article is transcribed and no reference is redistributed. Reference adaptation: none was needed and none was made; every reference inside the delivered unit resolves inside it, so no unresolved reference or citation is produced. Printed slips kept literal: no printed word, symbol, hypothesis or number of the counted statement or of its proof was altered. Layout and class: the article's own class is \texttt{amsart} with packages including geometry, amsmath, amsthm, amssymb, float, multirow, calligra, mathrsfs, subfig, inputenc, xcolor, hyperref, natbib, oubraces; the wrapper is the lane's pinned \texttt{amsart} editorial wrapper at 10pt on A4, which restates the theorem-like environments the retained text needs and adds the article's own macro definitions that the retained text needs (\texttt{\textbackslash bnabla}, \texttt{\textbackslash bx}, \texttt{\textbackslash calP}, \texttt{\textbackslash calT}, \texttt{\textbackslash jump}, \texttt{\textbackslash k}, \texttt{\textbackslash mean}, \texttt{\textbackslash mesh}, \texttt{\textbackslash ndgp}, \texttt{\textbackslash norm}, \texttt{\textbackslash ud}), with extra wrapper packages (bm, bbm, dsfont, xspace, cancel, mathtools). The delivered numbering is the unit's own. Assets: no retained span contains a figure environment, a graphics-inclusion command, a PDF-page inclusion, a table or a TikZ picture; no image file is copied into this package, the package declares no asset dependency and no image file is redistributed with it. Licence: version arXiv:2510.02094v1 of this article is distributed under the Creative Commons Attribution 4.0 International licence, as recorded in the arXiv licence metadata for this exact version and shown on the abstract page https://arxiv.org/abs/2510.02094v1. A full rights notice is delivered at licenses/arxiv-2510.02094-CC-BY-4.0.txt. Characters: every retained excerpt is pure ASCII, so the delivered engine, XeLaTeX with its default Latin Modern text font, reports no missing character.
  \begin{equation}\label{indef_est_important}
  \norm{ \sigma_\mesh^{-1/2}\mean{\mathcal{D}\nabla v}}{0,\Gamma_{\calP}}^2
  \le 
  16^{-1}\norm{\mathcal{D}^{1/2} \bnabla v}{0,\Omega}^2,
  \end{equation}

\begin{equation}
  s(w_h,v_h)=\sum_{K\in\mesh}\alpha\sum_{T\in\calT_K}\sum_{i=1}^{m_{k,d}}
  \big(\mathcal{D}_{\omega_K} h_T^{d-2}+\mu_T h_T^{d} \big)
  w_h(\bx_i^T)v_h(\bx_i^T) ,\label{eq20}
\end{equation}

\begin{lemma}
  \label{Lem:s}
There exists a constant $C_{\rm equiv}>0$, depending only on
$C_{\rm star}$, $C_{\rm sh}$, $k$, and the spatial dimension $d$, such
that for every $v_h\in W_\calT$,
\begin{align}
  \ndgp{ v_{h}}^{2}
   \le  C_{\rm equiv} \|\alpha^{-1/2} v_h \|_{s}^{2}.  \label{eq22}
\end{align}
\end{lemma}

\begin{lemma}\label{Lem:cont_twoscale}
Let $v_H,w_H\in V_\calP$, and set
$r_h^\pm:=\mathcal{E}^\pm(v_H)-\mathcal{E}^\pm(w_H)$ for brevity.
Define
\[
\calT_\k^\partial:=\{T\in \calT_\k:\ \exists \text{ face } f\subset \partial T\cap\partial K\}
\]
as the set of simplices $T$ in $\k$ touching the boundary of $\k$.  
Select $\alpha>0$ in \eqref{eq20} as
\begin{equation}\label{alpha_choice}
  \alpha|_\k:= \gamma \max\Big\{1,\frac{ H_{\k}}{8C_{\rm star}}\max_{T\in\calT_\k^\partial}\frac{|f|}{|T|}\Big\},\quad K\in\mesh,
\end{equation}
with $\gamma\ge 25C_{\rm equiv}$. Then
\begin{align}\label{cont_nonlinear}
a_{DG}(r_h^+, z_H) +s(r_h^-,z_H) 
\le \frac{7}{4}\Big(\ndgp{r_h^+}^2+\norm{r_h^-}{s}^2\Big)^{1/2} 
\Big(\ndgp{z_H}^2+\norm{z_H}{s}^2\Big)^{1/2}, 
\end{align}
for any $z_H\in V_\mesh$.
\end{lemma}

\begin{proof}
From standard estimates, we obtain
\begin{equation}\label{cont_one}
a_{DG}(r_h^+, z_H) \le \ndgp{r_h^+}\ndgp{z_H}
-	\int_{\Gamma_{\calP}}\mean{\mathcal{D}\nabla r_h^+}\cdot\jump{z_H}\ud s
-	\int_{\Gamma_{\calP}}\mean{\mathcal{D}\nabla z_H}\cdot\jump{r_h^+}\ud s.
\end{equation}

We now estimate the second and third terms on the right-hand side.  
For the last term, applying Cauchy--Schwarz together with
\eqref{indef_est_important} gives
\begin{equation}\label{indef_easy}
\Big|\int_{\Gamma_{\calP}}\mean{\mathcal{D}\nabla z_H}\cdot\jump{r_h^+}\ud s\Big|
\le 
4^{-1}\norm{\mathcal{D}^{1/2}\bnabla z_H}{0,\Omega}
\norm{\sqrt{\sigma_\calP}\jump{r_h^+}}{0,\Gamma_{\calP}}.
\end{equation}

For the remaining term, we decompose as
\begin{equation}\label{indef_cool}
\int_{\Gamma_{\calP}}\mean{\mathcal{D}\nabla r_h^+}\cdot\jump{z_H}\ud s
= \int_{\Gamma_{\calP}}\mean{\mathcal{D}\nabla (v_H-w_H)}\cdot\jump{z_H}\ud s
-\int_{\Gamma_{\calP}}\mean{\mathcal{D}\nabla r_h^-}\cdot\jump{z_H}\ud s
=:(I)-(II),
\end{equation}
since $r_h^++r_h^-=v_H-w_H$.

For $(I)$, Cauchy--Schwarz and \eqref{indef_est_important} yield
\[
\begin{aligned}
|(I)|
\le&\ 
4^{-1}\norm{\mathcal{D}^{1/2}\bnabla (v_H-w_H)}{0,\Omega}
\norm{\sqrt{\sigma_\calP}\jump{z_H}}{0,\Gamma_{\calP}}\\
\le&\
4^{-1}\big(\norm{\mathcal{D}^{1/2}\bnabla r_h^+}{0,\Omega}
+\norm{\mathcal{D}^{1/2}\bnabla r_h^-}{0,\Omega}\big)
\norm{\sqrt{\sigma_\calP}\jump{z_H}}{0,\Gamma_{\calP}}.
\end{aligned}
\]

For $(II)$, applying a trace--inverse inequality to every
$T\in\calT_\k^\partial$, $\k\in\mesh$, gives
\[
\begin{aligned}
|(II)|\le&\
\sum_{\k\in\mesh}\sum_{\substack{f\subset \partial T\cap\partial K\\ T\in \mathcal{T}_\k^\partial}}
\norm{\sigma_\calP^{-1/2}\mathcal{D}\nabla r_h^-|_T}{0,f}
\norm{\sqrt{\sigma_\calP}\jump{z_H}}{0,f}\\
\le&\
\bigg(\sum_{\k\in\mesh}\sum_{\substack{T\in \mathcal{T}_\k^\partial}}
\frac{\delta_{\k}|f|}{8C_{\rm star}|T|}
\min_{*\in\{+,-\}}H_{\k_*}\delta_{\k_*}^{-1}
\norm{\mathcal{D}^{1/2}\nabla r_h^-}{0,T}^2\bigg)^{1/2}
\norm{\sqrt{\sigma_\calP}\jump{z_H}}{0,\Gamma_{\calP}}.
\end{aligned}
\]

Since $\min_{*\in\{+,-\}}H_{\k_*}\delta_{\k_*}^{-1}\le H_\k\delta_{\k}^{-1}$ (as one of the $\k_*$ is $\k$ itself), and
\[
\alpha|_\k= \gamma \max\Big\{1,\frac{ H_{\k}}{8C_{\rm star}}\max_{T\in\calT_\k^\partial}\frac{|f|}{|T|}\Big\},
\]
we deduce
\begin{align}\label{cool_minus}
|(II)|
\le&\ \gamma^{-1/2}\bigg(\sum_{\k\in\mesh}\sum_{T\in \mathcal{T}_\k^\partial}
\norm{\sqrt{\alpha}\mathcal{D}^{1/2}\nabla r_h^-}{0,T}^2\bigg)^{1/2}
\norm{\sqrt{\sigma_\calP}\jump{z_H}}{0,\Gamma_{\calP}}\\
\le&\
\gamma^{-1/2}\norm{\sqrt{\alpha}\mathcal{D}^{1/2}\bnabla r_h^-}{0,\Omega}
\norm{\sqrt{\sigma_\calP}\jump{z_H}}{0,\Gamma_{\calP}}.\nonumber
\end{align}

Returning to \eqref{indef_cool}, the above bounds and Lemma
\ref{Lem:s}, with $r_h^-\in W_\calT$ and $\sqrt{\alpha}r_h^-\in
W_\calT$, imply
\[
\begin{aligned}
\Big|	\int_{\Gamma_{\calP}}\mean{\mathcal{D}\nabla r_h^+}\cdot\jump{z_H}\ud s\Big|
\le&\  4^{-1}\Big(\norm{\mathcal{D}^{1/2}\bnabla r_h^+}{0,\Omega}
+\norm{\mathcal{D}^{1/2}\bnabla r_h^-}{0,\Omega}
+4\gamma^{-1/2}\norm{\sqrt{\alpha}\mathcal{D}^{1/2}\bnabla r_h^-}{0,\Omega}\Big)
\\
&
\qquad \times\norm{\sqrt{\sigma_\calP}\jump{z_H}}{0,\Gamma_{\calP}}\\
\le&\  \tfrac{\sqrt{2}}{4}\Big(\norm{ \mathcal{D}^{1/2}\bnabla r_h^+}{0,\Omega}^2
+ 25 C_{\rm equiv}\gamma^{-1}\norm{ r_h^-}{s}^2\Big)^{1/2}
\norm{\sqrt{\sigma_\calP}\jump{z_H}}{0,\Gamma_{\calP}}.
\end{aligned}
\]

Since $\gamma\ge 25C_{\rm equiv}$, we conclude
\begin{equation}\label{indef_cool_too}
\Big|\int_{\Gamma_{\calP}}\mean{\mathcal{D}\nabla r_h^+}\cdot\jump{z_H}\ud s\Big|
\le 
\tfrac{\sqrt{2}}{4}\big(\norm{\mathcal{D}^{1/2}\bnabla r_h^+}{0,\Omega}^2
+\norm{r_h^-}{s}^2\big)^{1/2}
\norm{\sqrt{\sigma_\calP}\jump{z_H}}{0,\Gamma_{\calP}}.
\end{equation}

Finally, substituting \eqref{indef_easy} and \eqref{indef_cool_too}
into \eqref{cont_one}, and using $\sqrt{2}\le 2$, yields
\eqref{cont_nonlinear}.
\end{proof}

\end{document}
